\begin{abstract}
LLM judges are widely used to evaluate agentic tool-calling systems, yet their reliability on structured, dependency-driven workflows remains largely unexamined. We present AgentJudgeBench, the first benchmark to systematically study LLM-as-a-judge reliability for agentic tool-calling over workflow DAGs, as distinct from the broader LLM-as-a-judge task of open-ended text or preference evaluation. The benchmark comprises 3,808 instances spanning six DAG topologies and three difficulty tiers, evaluated with five generators (3B--70B open-weight models and GPT-5.4) and six judges (20B to frontier scale) under paired with- and without-ground-truth conditions. Judge alignment degrades monotonically with task difficulty, 1.5$\times$ faster without ground truth, and on hard queries without ground truth all six judges converge to a narrow 77--82\% band regardless of scale, revealing a structural ceiling, driven primarily by task difficulty though its height is partly prompt-dependent for weaker generators, that model capacity alone cannot overcome. Ground-truth exposure is not uniformly beneficial: it reduces alignment for GPT-5.4 (1.5 pp) and Gemini-2.5-Pro (3.9 pp), consistent with over-anchoring. Among mitigation strategies, chain-of-thought reasoning and judge temperature both have negligible effect, while structured evaluation rubrics improve alignment by up to 6.5 pp but do not generalise uniformly across judge-generator pairs. With ground truth, QwQ-32B best matches the programmatic reference, while a human validation study identifies GPT-OSS-120B as the most human-aligned judge; without it, frontier judges lead only marginally within the shared ceiling. These results expose fundamental limitations of current LLM judges and yield practical guidelines for reliable evaluation in agentic systems.
\end{abstract}